\subsection*{Appendix E: Causal-Study Design Brief}

\begin{tcolorbox}[colback=white,colframe=BayesSecondary!40,arc=2pt,boxrule=0.6pt,title={\small\textbf{The Causal-Study Design Brief}}]
\footnotesize\ttfamily
DECISION THE STUDY SHOULD INFORM: \\
ELIGIBLE POPULATION: \\
INTERVENTION AND COMPARISON CONDITIONS: \\
UNIT OF ASSIGNMENT AND UNIT OF ANALYSIS: \\
WHY THIS COMPARISON IS ETHICAL AND FEASIBLE: \\
PRIMARY OUTCOME AND HOW ITS VALIDITY WILL BE CHECKED: \\
PRIMARY ESTIMAND (E.G., OFFER EFFECT): \\
ASSIGNMENT PROCESS / SELECTION RISKS: \\
BASELINE DATA: \\
FOLLOW-UP HORIZON: \\
ANTICIPATED SAMPLE SIZE AND PRECISION LIMITS: \\
NONCOMPLIANCE / ATTRITION / MISSINGNESS PLAN: \\
POTENTIAL CONTAMINATION AND SPILLOVERS: \\
ANALYSIS PLAN AND IMPORTANT ASSUMPTIONS: \\
MINIMUM PRACTICALLY VALUABLE EFFECT: \\
BALANCING MEASURES AND SAFETY RULES: \\
HOW NEGATIVE OR INCONCLUSIVE FINDINGS WILL BE REPORTED:
\end{tcolorbox}

Do not use the mere completion of this brief to claim experimental rigor. It is a tool for revealing where rigor is not yet feasible and for deciding what less ambitious inference can still be useful.

\subsection*{Appendix F: A Practical Measurement Failure Table}

\begin{table*}[t]
\centering
\small
\caption{Practical Measurement Failures and Corrective Actions}
\begin{tabularx}{\textwidth}{p{0.28\textwidth}p{0.32\textwidth}X}
\toprule
\textbf{Symptom} & \textbf{Plausible Problem} & \textbf{Founder Response} \\
\midrule
Pass rates increase after a rubric revision & Changed assessment standard & Bridge or recalibrate before claiming progress \\
High marks but poor unseen performance & Rehearsal or answer dependence & Use unfamiliar tasks and an independence check \\
High satisfaction, low mastery & Experience and learning diverge & Repair instruction; measure both separately \\
Only strongest learners finish & Attrition/selection & Report all starters and investigate recovery paths \\
Reviewers disagree sharply & Unclear scoring or task construct & Calibrate and revise anchors \\
Great project demos but no authorship evidence & Assistance concealed & Record process, tools, and defense \\
Training completed, workplace unchanged & No transfer conditions & Work with sponsor on application permissions and workflow \\
Hours supposedly saved, P\&L unchanged & Capacity value not realized cash & Track genuine productive reallocation or cost effects \\
Improvement concentrated in 1 subgroup & Heterogeneous fit or support & Examine prerequisites and differentiated pathways \\
Dashboard green, complaints rising & Metric misses customer friction & Add meaningful balancing measures \\
Results cannot be reproduced next cohort & Cohort mix, faculty, task or operations changed & Document and compare conditions \\
Marketing claims stronger than available data & Inference inflation & Narrow and correct public statement \\
\bottomrule
\end{tabularx}
\end{table*}

\subsection*{Appendix G: The Five-Metric Early Educational Dashboard}

For a very small launch, avoid excessive reporting. A compact internal dashboard can begin with: \textbf{(1) entry and fit}, \textbf{(2) independent target-task evidence}, \textbf{(3) completion and attrition}, \textbf{(4) learner feedback and service reliability}, and \textbf{(5) full economic contribution including normalized founder hours}. Add delayed retention and external transfer as the program develops, clearly distinguishing unmeasured from 0.

The 5 metrics are not universally optimal and not substitutes for safeguarding and privacy. Their value is that they prevent a new institute from equating enrollment, satisfaction, or financial revenue with genuine transformation.
